\section{Elliptic period squares}\label{sec:general-period-square}

The distinction between existence and uniqueness becomes transparent in
de Rham cohomology. A first-order identity selects a line in a
two-dimensional algebraic space. Complex multiplication supplies that
line; Chudnovsky's theorem makes it unique. The unproved converse asks
whether the complex period pairing can select it in any other way.

\subsection{The determinant and the first jet}
Put $K=\Qbar$ and $\Pi=2\pi i$. Let $E\to X$ be a smooth elliptic
family over a smooth algebraic curve over $K$, and write
\[
 \mathcal H=H^1_{\mathrm{dR}}(E/X),\qquad
 \mathcal L=f_*\Omega^1_{E/X}\subset\mathcal H.
\]
Here $f:E\to X$ is the structure morphism. Choose a local algebraic
frame $(\lambda,\mu)$ of $\mathcal H$, with $\lambda$ spanning
$\mathcal L$ and normalized cup product $\langle\lambda,\mu\rangle=1$.
At an algebraic fiber $P\in X(K)$, viewed over $\C$, the comparison
isomorphism respects the determinant:
\[
\begin{tikzcd}[column sep=large]
 \bigwedge^2 H^1_{\mathrm{dR}}(E_P)\otimes_K\C
   \arrow[r,"\bigwedge^2\mathrm{comp}"]
   \arrow[d,"{\langle\ ,\ \rangle}"']
 &\bigwedge^2 H^1_B(E_P,\Q)\otimes_\Q\C
   \arrow[d,"\mathrm{ev}_{\gamma_1\wedge\gamma_2}"]\\
 \C\arrow[r,"\times\Pi"']&\C .
\end{tikzcd}
\]
The oriented integral basis $(\gamma_1,\gamma_2)$ is chosen compatibly
with this normalization. Thus the Legendre relation is
\begin{equation}\label{eq:general-legendre}
 \omega_1h_2-\omega_2h_1=\Pi,
 \qquad \omega_j=\int_{\gamma_j}\lambda,\quad
 h_j=\int_{\gamma_j}\mu.
\end{equation}
The determinant is a Tate period; an individual product $\omega h$
need not be one.

Let $D$ be an algebraic derivation, $\gamma$ a nonzero locally flat
integral cycle, and $k$ a nonzero algebraic function. Suppose the
analytic function under consideration has the realization
\begin{equation}\label{eq:general-square}
 H=k\frac{\omega^2}{\Pi^2},\qquad
 \omega=\int_\gamma\lambda,\quad h=\int_\gamma\mu.
\end{equation}
Write
\begin{equation}\label{eq:general-connection}
 \nabla_D\lambda=p\lambda+q\mu,\qquad u=D\log k+2p.
\end{equation}
A specialization $P\in X(K)$ is \emph{admissible} if all these data
are regular there and $k(P)q(P)\ne0$. The condition $q(P)\ne0$ says
that the Kodaira--Spencer map evaluated on $D$ is nonzero. At such a
point the coefficient map
\begin{equation}\label{eq:general-coefficient-class}
 \begin{gathered}
 T_P:K^2\xrightarrow{\ \sim\ }\mathcal H_P,\qquad
 (A,B)\longmapsto\nu_{A,B},\\
 \nu_{A,B}=(A+B D\log k)\lambda+2B\nabla_D\lambda .
 \end{gathered}
\end{equation}
is an isomorphism. Differentiating the period pairing gives
\begin{equation}\label{eq:general-jet}
 (A+BD)H(P)
   =\frac{k\omega}{\Pi^2}\int_\gamma\nu_{A,B}
   =\frac{k}{\Pi^2}\bigl((A+Bu)\omega^2+2Bq\omega h\bigr).
\end{equation}
All quantities in a fiber formula are henceforth evaluated at $P$.
This construction is intrinsic: replacing $\lambda$ by $a\lambda$
and $k$ by $k/a^2$ replaces $\nu_{A,B}$ by $a\nu_{A,B}$, leaving the
period pairing and the selected line unchanged.

\subsection{The transcendence input}
\begin{theorem}[Chudnovsky]\label{thm:B-chudnovsky}
Let $E/K$ be an elliptic curve, $(\lambda,\mu)$ an algebraic de Rham
basis with $\lambda$ holomorphic, and $\gamma$ a nonzero integral
cycle. Then
\begin{equation}\label{eq:chudnovsky}
 R=\frac h\omega,\qquad S=\frac\pi\omega
 \quad\text{are algebraically independent over }K.
\end{equation}
\end{theorem}
We specify the source and normalization, because the power of
$\omega$ in this assertion is essential. For a Weierstrass equation
$y^2=4X^3-g_2X-g_3$ over $K$, take $\lambda=dX/y$ and
$\mu=X\,dX/y$. Uniformization gives $h_j=-\eta_j$, where $\eta_j$
is the zeta quasi-period of the \emph{full} lattice period $\omega_j$.
Consequently
\begin{equation}\label{eq:B-legendre-convention}
 \omega_2\eta_1-\omega_1\eta_2=2\pi i.
\end{equation}
Corollary~4 of \cite{Waldschmidt}, equivalently Corollary~31 of
\cite[p.~166]{WaldschmidtSurvey}, gives precisely
\eqref{eq:chudnovsky}. To see its origin in the elliptic-logarithm
form of Chudnovsky's theorem, complete a primitive $\omega_1$ to a
period basis and take the algebraic two-torsion point with logarithm
$v=\omega_2/2$. Then $\wp(v)\in K$, $v$ and $\omega_1$ are
$\Q$-linearly independent, and
\[
 \zeta(v)-\frac{\eta_1}{\omega_1}v
   =\frac12\left(\eta_2-\frac{\eta_1\omega_2}{\omega_1}\right)
   =-\frac{\pi i}{\omega_1}.
\]
This is the second independent quantity in the cited theorem.
Replacing the Weierstrass basis by any algebraic basis as above
changes $R$ by an invertible algebraic affine transformation and $S$
by a nonzero algebraic factor. Multiplying the cycle by a nonzero
integer only rescales $S$. These operations preserve algebraic
independence. In particular,
\begin{equation}\label{eq:general-ratio-transcendence}
                         h/\omega\notin K.
\end{equation}

\subsection{The identity line}
\begin{theorem}\label{thm:general-identity-line}
At an admissible specialization let
\[
 \mathcal R_P=\{(A,B)\in K^2:(A+BD)H(P)\in K/\Pi\}.
\]
Then $\dim_K\mathcal R_P\le1$, and a nonzero coefficient pair never
has value zero. If $E_P$ has CM, $\mathcal R_P$ is exactly the
inverse image under $T_P$ of the CM eigenline complementary to
$\mathcal L_P$. More explicitly, if that eigenline is
$K(\mu-r\lambda)$, then
\begin{equation}\label{eq:general-cm-slope}
 \mathcal R_P=K\bigl(-(u+2qr),1\bigr).
\end{equation}
Thus a rational coefficient pair exists exactly when
$u+2qr\in\Q$, and its primitive integer representative is unique
up to sign.

If $\Pi,\omega^2,\omega h$ are $K$-linearly independent at every
admissible non-CM fiber in question, then
$\mathcal R_P\ne0$ is equivalent to CM.
\end{theorem}
\begin{proof}
The holomorphic period $\omega$ is nonzero: integration identifies
integral homology with a lattice in $\C$. By
\eqref{eq:general-ratio-transcendence}, evaluation on $\gamma$ is
injective on the two-dimensional $K$-space $\mathcal H_P$.
Equation~\eqref{eq:general-jet} and the isomorphism $T_P$ therefore
make $(A,B)\mapsto(A+BD)H(P)$ injective. Its inverse image of the
one-dimensional space $K/\Pi$ has dimension at most one.

Suppose $E_P$ has CM. A nonreal endomorphism $\varphi$ has distinct
algebraic eigenvalues $a,b$ on $\mathcal H_P$, with $a$ on the
holomorphic line. Its other eigenline has the form
$K\nu$, where $\nu=\mu-r\lambda$ and $r\in K$. The cycles
$\gamma,\varphi_*\gamma$ are independent over $\Q$, since a
rational dependence would make the nonreal eigenvalue $a$ rational.
If their determinant in an oriented integral basis is
$m\in\Z\setminus\{0\}$, compatibility of the endomorphism with
the comparison isomorphism gives
\begin{equation}\label{eq:general-cm-determinant}
 (b-a)\omega\int_\gamma\nu=m\Pi,\qquad
 c_P:=\frac{\omega(h-r\omega)}{\Pi}
       =\frac{m}{b-a}\in K^\times.
\end{equation}
Hence the pair in \eqref{eq:general-cm-slope} has value
$2kq c_P/\Pi\ne0$. Uniqueness identifies its line with all of
$\mathcal R_P$. The rational and primitive assertions follow from
its representative with second coordinate one.

Finally a scalar identity is equivalent to
\begin{equation}\label{eq:general-period-relation}
 k(A+Bu)\omega^2+2kBq\omega h=\alpha\Pi.
\end{equation}
The two coefficients on the left cannot both vanish for a nonzero
pair. The asserted independence therefore excludes a non-CM fiber.
\end{proof}

This last independence hypothesis is equivalent, at the specified
fiber, to exclusion of \emph{all algebraic} coefficient identities:
the triangular coefficient transformation in
\eqref{eq:general-period-relation} is invertible. For the original
problem only rational coefficient directions need be excluded. The
distinction will be retained in the converse criteria below.

\begin{corollary}[CM period independence]\label{prop:B-CM-independent}
If $E/K$ has CM, every nonzero holomorphic period $\omega$ and
$\pi$ are algebraically independent. For the corresponding
second-kind period there are $c,d\in K$, $d\ne0$, such that
\begin{equation}\label{eq:B-CM-quasiperiod}
                        h=c\omega+d\frac\pi\omega.
\end{equation}
\end{corollary}
\begin{proof}
Equation~\eqref{eq:general-cm-determinant} gives this formula with
$c=r$ and $d=2ic_P$. Consequently
\[
 K(\omega,\pi)=K(R,S),\qquad
 \omega=\frac{dS}{R-c},\quad \pi=\frac{dS^2}{R-c}.
\]
Chudnovsky's theorem makes this field have transcendence degree two.
This proves the exact CM consequence also stated in
\cite[Corollary~6]{Waldschmidt} and
\cite[Corollary~32]{WaldschmidtSurvey}; the algebraic discriminant
factor in the former normalization does not affect the assertion.
\end{proof}

\subsection{The four standard realizations}\label{sec:families}
The general construction becomes uniform after the quadratic change
\begin{equation}\label{eq:clausen}
 x=\frac{1-\sqrt{1-z}}2,\quad z=4x(1-x),\qquad
 F_s(z)=f_s(x)^2,\quad f_s(x)={}_2F_1(s,1-s;1;x).
\end{equation}
The square root is nonnegative on $[-1,1]$. The equality follows
from Clausen's identity and the Gauss quadratic transformation,
first as analytic germs and then along this real interval.
The models $E_{s,x}$ and their de Rham calculation are given in
Appendix~\ref{rev:period-models}. With
$\lambda=dX/y$ and $\mu=X\,dX/y$ they give the following common
normalization.

\begin{lemma}\label{lem:period}
The cycle analytic at the nodal fiber $x=0$ has an orientation for
which
\begin{equation}\label{eq:normalization}
 \omega_s(x)=\frac{2\pi}{\sqrt6}f_s(x),\qquad
 F_s(z)=\frac{3\omega_s(x)^2}{2\pi^2}.
\end{equation}
For $z\in[-1,1)\setminus\{0\}$ the fiber is smooth and
\begin{equation}\label{eq:logderivative}
 \begin{gathered}
 \frac{\thetaop F_s}{F_s}
    =u_s+v_s\frac{h_s}{\omega_s},\\
 u_s=-\frac{x(1-x)}{6(1-2x)}\frac{\Delta_s'}{\Delta_s},\quad
 v_s=-\frac{\varepsilon_s}{3(1-2x)}\ne0,
 \end{gathered}
\end{equation}
where $\varepsilon_s=1$ for $s=1/6,1/4,1/3$ and
$\varepsilon_{1/2}=2$.
\end{lemma}
\begin{proof}
The Gauss--Manin connection in Appendix~\ref{rev:period-models}
shows that $\omega_s$ satisfies the Gauss equation
\[
 x(1-x)\omega_s''+(1-2x)\omega_s'-s(1-s)\omega_s=0.
\]
All four cubics specialize to $(X-1)(2X+1)^2$ at zero. The residue
at the node on the normalization is $\pm1/\sqrt{-6}$; the small
cycle therefore has limiting period $\pm2\pi/\sqrt6$.
The Gauss equation has a unique analytic solution with this chosen
positive value. This proves the first formula, and
\eqref{eq:clausen} proves the second. The same connection and
$\thetaop=x(1-x)(1-2x)^{-1}d/dx$ give
\eqref{eq:logderivative}. Euler's integral for $f_s$ is positive
for real $x<1$, so the period does not vanish on the interval in use.
\end{proof}
Thus $k=-6$, $D=\thetaop$, and $2q=v_s$ in the general theory.
In particular its nondegeneracy hypothesis holds at every standard
parameter below the positive endpoint.

\subsection{Value independence and uniqueness}\label{sec:unique}
\begin{theorem}\label{thm:B-value-independent}
For algebraic $z\in[-1,1)\setminus\{0\}$ the two values
$F_s(z),\thetaop F_s(z)$ are algebraically independent over $K$.
If $q_0=e^{2\pi i\tau}$ is the nome obtained by completing the
selected primitive cycle to an oriented period basis, then
$q_0,F_s(z),\thetaop F_s(z)$ are algebraically independent.
\end{theorem}
\begin{proof}
Put $F=F_s(z)$ and $G=\thetaop F_s(z)$. Equations
\eqref{eq:normalization} and \eqref{eq:logderivative} give
\[
 F=\frac3{2S^2},\qquad G=F(u_s+v_sR),\qquad
 K(F,G)=K(R,S^2).
\]
The last field has transcendence degree two by Chudnovsky's theorem.

For the nome statement use the full conclusion of Nesterenko's
theorem, in \cite[Theorem~1.1]{Fonseca}:
\[
 \operatorname{trdeg}_K K(q_0,E_2(\tau),E_4(\tau),E_6(\tau))\ge3.
\]
The period formulas in \cite[Proposition~4.7]{Fonseca}, in our
second-kind convention, are
\[
 E_2=-2FR,\qquad E_4=\frac{g_2}{3}F^2,\qquad E_6=g_3F^3.
\]
They place this field inside $K(q_0,F,G)$, a field with three
generators. Its transcendence degree is therefore three. If the
chosen cycle is an integer multiple of a primitive cycle, passing
to that primitive cycle changes these formulas only by nonzero
algebraic scaling factors and gives the same conclusion.
\end{proof}

\begin{corollary}\label{prop:B-relation-space}
At every such parameter the space
\[
 \{(A,B,\alpha)\in K^3:AF_s(z)+B\thetaop F_s(z)=\alpha/\pi\}
\]
has dimension at most one, and its nonzero elements have
$\alpha\ne0$. In particular there is at most one primitive integer
pair up to sign, and
\begin{equation}\label{eq:translog}
                    \thetaop F_s(z)/F_s(z)\notin K.
\end{equation}
\end{corollary}
\begin{proof}
Theorem~\ref{thm:general-identity-line} gives the assertion; the
multiplier is uniquely determined by the coefficient pair. Equivalently,
linear independence of $F_s$ and $\thetaop F_s$ makes projection onto
the $\alpha$ coordinate injective on the displayed relation space.
\end{proof}

\subsection{Ordinary convergence and the endpoint}\label{sec:endpoint}
Stirling's formula gives
\begin{equation}\label{eq:asymptotic}
 c_s(n)\sim
 \frac{n^{-3/2}}{\Gamma(1/2)\Gamma(s)\Gamma(1-s)}.
\end{equation}
For a nonzero pair $(A,B)$, the series is absolutely convergent for
$|z|<1$ and divergent for $|z|>1$. At $z=1$ it converges precisely
when $B=0$. At $z=-1$ the $A$ part converges absolutely, while the
$B$ part, if present, converges conditionally: the positive sequence
$nc_s(n)$ tends to zero and its successive ratio is
$1-1/(2n)+O(n^{-2})$. Abel's theorem identifies this sum with the
radial value used above.

\begin{proposition}
No nonzero algebraic coefficient pair gives an ordinarily convergent
value in $K/\pi$ at $z=1$.
\end{proposition}
\begin{proof}
Only $B=0$ is possible. The four smooth endpoint fibers have
$j$-invariants $1728,8000,54000,1728$, respectively, and all have
CM. One may verify this directly: $j=1728$ has an order-four
automorphism; $y^2=x^3+4x^2+2x$, of invariant $8000$, is
2-isogenous to an algebraically isomorphic curve; and the
invariant $54000$ occurs on a curve 2-isogenous to $y^2=x^3+1$.
The isogenies are recorded in Appendix~\ref{rev:period-models}.
Corollary~\ref{prop:B-CM-independent} makes $\omega_s(1/2)$ and
$\pi$ algebraically independent. But a nonzero identity at this
endpoint would make
$\pi F_s(1)=3\omega_s(1/2)^2/(2\pi)$ algebraic, a contradiction.
\end{proof}

The endpoint proposition and Corollary~\ref{prop:B-relation-space}
prove projective uniqueness throughout the ordinary convergence
range. They do not turn the independence of $h/\omega$ and
$\pi/\omega$ into independence of $\pi,\omega^2,\omega h$.
The latter is precisely the additional comparison required for CM
forcing.
